% !TeX root = ../../Book5.tex
% 纲要标题：### 34.6 四点标记有理曲线的稳定紧化
% 纲要位置：几何学计划纲要.md，第 6525 行。
% 正文状态：未撰写；本文件目前是框架。

% 纲要内容开始
% 
% 将一个紧化问题做完：主模型可在任意域上构造；复解析参数与长颈解释见 IV.6.5。先完成这个低维代表对象，再把一般稳定曲线模空间作为具有额外存在定理的研究出口。
% 
% * 有序标记、同构及交比归一化
% * $\mathcal M_{0,4}\simeq\mathbb P^1\setminus\{0,1,\infty\}$ 与光滑通用族
% * 四个截面的碰撞与三个吹起中心
% 
%   从 $\mathbb P^1_z\times\mathbb P^1_\lambda$ 及截面 $z=0,1,\infty,\lambda$ 出发，在 $(0,0)$、$(1,1)$、$(\infty,\infty)$ 吹起，取严格变换后的标记截面。
% 
% * 结点纤维、每个分支上的特殊点和稳定性
% * 三个边界点与标记集的三种二二分割
% * 平坦性、截面不交及稳定族的局部核验
% * 交比的族版本与代表对象的通用性
% 
%   除几何点分类外，还需在非约化测试基底上构造拉回对应，并验证局部平凡化的下降；不能仅凭交比给出几何点双射就宣布细模可表。
% 
% * $\overline{\mathcal M}_{0,4}\simeq\mathbb P^1$ 的固有性和粗细模辨别
% 
%   光滑四点有理曲线有四个标记点；边界纤维的每个分支恰有三个特殊点，逐支检查自同构平凡；置换标记后的商问题另行处理。
% 
% * 标记置换、商构造与更高亏格所需的新输入
% 
% ---
% 
% 纲要内容结束

\section{四点标记有理曲线的稳定紧化}
\label{g5:sec:3406}

待撰写

\subsection{有序标记、同构及交比归一化}
\label{g5:subsec:3406:01}
待撰写

\subsection{\texorpdfstring{\(\mathcal M_{0,4}\simeq\mathbb P^1\setminus\{0,1,\infty\}\) 与光滑通用族}{M\_\{0,4\}≃P¹∖\{0,1,∞\} 与光滑通用族}}
\label{g5:subsec:3406:02}
待撰写

\subsection{四个截面的碰撞与三个吹起中心}
\label{g5:subsec:3406:03}
待撰写

\subsection{结点纤维、每个分支上的特殊点和稳定性}
\label{g5:subsec:3406:04}
待撰写

\subsection{三个边界点与标记集的三种二二分割}
\label{g5:subsec:3406:05}
待撰写

\subsection{平坦性、截面不交及稳定族的局部核验}
\label{g5:subsec:3406:06}
待撰写

\subsection{交比的族版本与代表对象的通用性}
\label{g5:subsec:3406:07}
待撰写

\subsection{\texorpdfstring{\(\overline{\mathcal M}_{0,4}\simeq\mathbb P^1\) 的固有性和粗细模辨别}{\={M}\_\{0,4\}≃P¹ 的固有性和粗细模辨别}}
\label{g5:subsec:3406:08}
待撰写

\subsection{标记置换、商构造与更高亏格所需的新输入}
\label{g5:subsec:3406:09}
待撰写
% 纲要条目排为编号小节标题，并进入目录与书签；说明段仍保留于纲要注释中。
